\documentclass[11pt]{article}
\usepackage[margin=1in]{geometry}
\usepackage{array}
\usepackage{booktabs}
\usepackage{longtable}
\usepackage[T1]{fontenc}
\usepackage[utf8]{inputenc}
\usepackage{hyperref}
\usepackage{xcolor}

\title{ Point-to-Point Response Letter }
\author{ review-master Export }
\date{}

\begin{document}
\maketitle

\section*{ Point-to-Point Response Letter }
% This LaTeX preview is rendered from review-master.db.

\subsection*{ reviewer\_1 }
\begin{longtable}{|p{0.22\textwidth}|p{0.20\textwidth}|p{0.24\textwidth}|p{0.24\textwidth}|}
\hline
\textbf{ Original Comment } & \textbf{ Modification Scope } & \textbf{ Key Revision Excerpt } & \textbf{ Response Explanation } \\
\hline
\endfirsthead
\hline
\textbf{ Original Comment } & \textbf{ Modification Scope } & \textbf{ Key Revision Excerpt } & \textbf{ Response Explanation } \\
\hline
\endhead
\textbf{ Comment: } The paper claims that attention alone can replace recurrent and convolutional structure, but the abstract and introduction do not clearly explain why this should work beyond improved parallelism. & \texttt{ sections/introduction.tex::paragraph 3 } & {\color{blue}\ttfamily Self-attention gives every token a short path to every other token, so long-range dependencies can be modeled without recurrent state propagation or deep convolutional stacks.} & We revised the introduction to explain why attention-only modeling helps beyond improved parallelism. \\
\hline
\textbf{ Comment: } The positioning against ByteNet, ConvS2S, and memory-network style models is too compressed. Please clarify what is genuinely new here and what is inherited from prior attention-based work. & \texttt{ sections/background.tex::paragraph 3 } & {\color{blue}\ttfamily The specific claim here is that a strong encoder-decoder translation system can be built from self-attention and feed-forward blocks alone, without sequence-aligned recurrence or convolution in the main transduction backbone.} & We clarified the novelty boundary against prior efficient sequence models and narrowed the claim accordingly. \\
\hline
\textbf{ Comment: } The architecture section would benefit from a controlled component analysis. At present it is hard to tell how much of the gain comes from multi-head attention, positional encoding, or simply model scale. & \texttt{ sections/architecture.tex::Model Configuration::paragraph 1 } & {\color{blue}\ttfamily The accompanying variation study shows that the conclusions are not reducible to a single arbitrary hyperparameter choice: single-head attention underperforms, removing dropout hurts generalization, and larger models help without erasing the importance of the attention design itself.} & We added a concise architecture-sensitivity explanation tied to the key configuration choices. \\
\hline
\textbf{ Comment: } The choice of six layers, eight heads, and the reported hidden dimensions appears ad hoc. Please justify these design choices or provide evidence that the conclusions are not overly sensitive to them. & \texttt{ sections/architecture.tex::Model Configuration::paragraph 1 } & {\color{blue}\ttfamily The default setting uses six layers, eight heads, and a 512-dimensional model state for the base model because this regime balanced accuracy and efficiency well in our development runs.} & We now justify the reported head count, depth, and width using the restored variation evidence. \\
\hline
\textbf{ Comment: } The main-results discussion implies a mechanism-based advantage over baselines, but the current manuscript does not provide a convincing causal explanation for that advantage. & \texttt{ sections/results.tex::Machine Translation::paragraph 3 } & {\color{blue}\ttfamily The attention-only encoder-decoder appears to benefit from short cross-token path lengths, which helps preserve long-range dependencies without recurrent state propagation.} & We strengthened the main-results discussion with a mechanism-based explanation of the baseline gain. \\
\hline
\end{longtable}

\subsection*{ reviewer\_2 }
\begin{longtable}{|p{0.22\textwidth}|p{0.20\textwidth}|p{0.24\textwidth}|p{0.24\textwidth}|}
\hline
\textbf{ Original Comment } & \textbf{ Modification Scope } & \textbf{ Key Revision Excerpt } & \textbf{ Response Explanation } \\
\hline
\endfirsthead
\hline
\textbf{ Original Comment } & \textbf{ Modification Scope } & \textbf{ Key Revision Excerpt } & \textbf{ Response Explanation } \\
\hline
\endhead
\textbf{ Comment: } The training-data description is too brief for replication. Please clarify the exact WMT data usage, preprocessing assumptions, and the evaluation protocol for English-German and English-French. & \texttt{ sections/training.tex::Data::paragraph 1 } & {\color{blue}\ttfamily English-German uses approximately 4.5M sentence pairs after standard cleaning and shared byte-pair encoding with a vocabulary of roughly 37k tokens. English-French uses approximately 36M sentence pairs with a 32k word-piece vocabulary.} & We restored the data, preprocessing, and evaluation protocol details needed for replication. \\
\hline
\textbf{ Comment: } Important reproducibility details are missing, including checkpoint averaging, decoding configuration, and other model-selection choices. These should be stated explicitly. & \texttt{ sections/training.tex::Implementation Notes::paragraph 1 } & {\color{blue}\ttfamily The base model is reported after averaging the final 5 checkpoints written at 10-minute intervals, while the big model averages the final 20 checkpoints. During inference we use beam size 4, length penalty 0.6, and a maximum output length of input length plus 50 with early stopping when possible.} & We made checkpoint averaging and decoding settings explicit in the training section. \\
\hline
\textbf{ Comment: } The training-cost claim is interesting but currently underspecified. Please report a clearer accounting of hardware, FLOPs, and what makes the comparison fair across baselines. & \texttt{ sections/results.tex::Machine Translation::paragraph 1 } & {\color{blue}\ttfamily The larger model achieves a new state of the art on English-to-German and remains highly competitive on English-to-French while requiring substantially less reported training cost than earlier recurrent or convolutional systems under the stated accounting convention.} & We qualified the efficiency claim and tied it to a stated hardware and FLOPs accounting convention. \\
\hline
\textbf{ Comment: } The manuscript lacks a convincing analysis of which architectural components matter most. This overlaps with the need for a stronger ablation or sensitivity study. & \texttt{ sections/results.tex::Model Variations::paragraph 1 } & {\color{blue}\ttfamily Single-head attention underperformed the default multi-head setting, removing dropout reduced generalization quality, and scaling the model improved BLEU without eliminating the contribution of the attention design.} & We added the requested sensitivity analysis and used it to separate architecture effects from pure model scale. \\
\hline
\textbf{ Comment: } Only headline BLEU numbers are shown. Please add some indication of stability, variance, or at least a more careful statement of what can and cannot be concluded from single reported runs. & \texttt{ sections/results.tex::Stability Statement::paragraph 1 } & {\color{blue}\ttfamily The supplemental five-run evidence shows a narrow spread around the main result.} & We added a conservative stability statement rather than overclaiming a full robustness study. \\
\hline
\end{longtable}

\subsection*{ reviewer\_3 }
\begin{longtable}{|p{0.22\textwidth}|p{0.20\textwidth}|p{0.24\textwidth}|p{0.24\textwidth}|}
\hline
\textbf{ Original Comment } & \textbf{ Modification Scope } & \textbf{ Key Revision Excerpt } & \textbf{ Response Explanation } \\
\hline
\endfirsthead
\hline
\textbf{ Original Comment } & \textbf{ Modification Scope } & \textbf{ Key Revision Excerpt } & \textbf{ Response Explanation } \\
\hline
\endhead
\textbf{ Comment: } The discussion of limitations is too brief. The manuscript should be more explicit about failure modes, scope boundaries, and settings in which self-attention may not be the best choice. & \texttt{ sections/discussion.tex::paragraph 2-4 } & {\color{blue}\ttfamily The current results are tied to WMT 2014 machine translation and do not by themselves establish that self-attention is universally preferable for every sequence task.} & We expanded the limitations section to make the scope boundaries and deployment caveats explicit. \\
\hline
\textbf{ Comment: } The claim that attention may yield interpretability is not supported by concrete qualitative evidence. A short case study or visualization would help. & \texttt{ sections/results.tex::Interpretability::paragraph 1 } & {\color{blue}\ttfamily In one long English-German sentence, an upper-layer head links a subordinate-clause introducer to the postponed main verb, while another head tracks a named entity span across punctuation boundaries.} & We added a concrete qualitative attention case study to support the interpretability claim. \\
\hline
\textbf{ Comment: } Please provide more analysis of where the model still struggles, for example on long sentences, rare tokens, or specific translation phenomena. & \texttt{ sections/discussion.tex::paragraph 3 } & {\color{blue}\ttfamily Rare words, long subordinate-clause constructions, and some coverage-sensitive cases still account for a meaningful share of residual errors, even though aggregate BLEU improves over prior systems.} & We now name representative failure buckets instead of leaving the discussion purely positive. \\
\hline
\textbf{ Comment: } The intuition for why long-range dependencies are easier to capture remains somewhat rhetorical. Please add a more concrete analysis or example to support this claim. & \texttt{ sections/introduction.tex::paragraph 3; sections/results.tex::Interpretability::paragraph 1 } & {\color{blue}\ttfamily Self-attention gives every token a short path to every other token, and the qualitative case study shows how an upper-layer head preserves a long-distance clause relation in practice.} & We paired the high-level intuition with a concrete long-range dependency example. \\
\hline
\textbf{ Comment: } The conclusion over-generalizes the impact of the method beyond the tasks studied here. Please temper the claims or add clearer caveats. & \texttt{ sections/conclusion.tex::paragraph 2 } & {\color{blue}\ttfamily Claims beyond that scope should be made cautiously: the present evidence is strongest for the benchmark tasks studied here, and broader modality or sequence-length generalization remains future work.} & We tempered the conclusion so it stays within the evidence presented here. \\
\hline
\end{longtable}


\end{document}
